\scanpage{7}
\mspage{3}
à savoir le système projectif des $B_{K_i}$, ou des $\Pi_{K_i}$, pour $K_i$ sous-corps de $K$ de type fini sur le corps premier, et à un syst. $(K,\bar K)$ le système projectif des $E_{\bar K_i/K_i}$, où $\bar K_i$ est la clôture alg. sép. de $K_i$ dans $\bar K$. On aura
\[
(4)\qquad\left\{\begin{aligned}
\Pi_K&\cong\varprojlim\Pi_{K_i}\\
B_K&\cong\varprojlim B_{K_i}\\
E_{\bar K/K}&\xrightarrow{\sim}\varprojlim E_{\bar K_i/K_{\unsure{i}}}
\end{aligned}\right.
\]
i.e. on reconstitue les objets $B_K$, $\Pi_K$, $E_{\bar K/K}$ à partir du syst. proj. correspondant -- mais l'inverse n'est pas vrai.

En fait, comme le foncteur $\operatorname{Ind}(\text{corps.t.f})\longrightarrow\text{corps}$ de la catégorie des syst. inductifs de corps de type fini, vers celle des corps, est une équivalence de catégories (pour des raisons triviales), il s'ensuit que le foncteur $K\longrightarrow B_K$, ou $K\unsure{-|\longrightarrow}\Pi_K$, ou $(K,\bar K)\longrightarrow E_{\bar K/K}$, \unsure{étendu|étant}
\margin{6}
